\documentclass[dvipdfmx,a4paper,10pt]{jsarticle}

\usepackage[top=20mm,bottom=22mm,left=20mm,right=20mm]{geometry}
\usepackage{amsmath,amssymb,amsfonts}
\usepackage{bm}
\usepackage{booktabs}
\usepackage{array}
\usepackage[haranoaji]{pxchfon}

\setlength{\parindent}{1em}
\setlength{\parskip}{0.25em}
\renewcommand{\arraystretch}{1.15}

\newcommand{\ii}{\mathrm{i}}
\newcommand{\tr}{\operatorname{Tr}}
\newcommand{\E}{\mathbb{E}}
\newcommand{\Var}{\operatorname{Var}}
\newcommand{\Hc}{\mathcal{H}}
\newcommand{\Dc}{\mathcal{D}}
\newcommand{\Cc}{\mathcal{C}}
\newcommand{\ket}[1]{\lvert #1\rangle}
\newcommand{\bra}[1]{\langle #1\rvert}

\title{Fock状態から導く有限Lamb--Dicke MSゲートの\\
有効誤差生成子}
\author{平井 希空}
\date{\today}

\begin{document}

\maketitle

\begin{abstract}
有限Lamb--Dicke（LD）補正を含むM{\o}lmer--S{\o}rensen（MS）ゲートに
ついて，数値的に観測された主要な三つの誤差成分をFock状態から解析的に
導出する．有限LD補正はFock数に依存する幾何学的位相と不完全な位相空間
軌道閉鎖を生じさせる．その熱平均を取ることにより，有効誤差生成子は
\[
 \mathcal{K}_{\mathrm{err}}
 \simeq
 h_{XX}\Hc[XX]
 +\gamma_{XX}\Dc[XX]
 +\gamma_{\mathrm{col}}\Dc[IX+XI]
\]
と表される．ここで三係数はそれぞれ，Fock依存回転角の平均，その熱分散，
および残留スピン振動エンタングルメントに対応する．一点較正後の第二Magnus予測と
第三Magnus予測は，$\bar n\leq4$の範囲でfull numerical evolutionから
得た係数を約$3.1\%$以内で再現する．
\end{abstract}

\section{有効誤差生成子}

本稿では
\begin{align}
 \Hc[P](\rho)&=-\ii[P,\rho],\\
 \Dc[L](\rho)&=
 L\rho L^\dagger-\frac{1}{2}\{L^\dagger L,\rho\}
\end{align}
と定義する．導出する三成分は
\begin{equation}
 \boxed{
 \mathcal{K}_{\mathrm{err}}
 \simeq
 h_{XX}\Hc[XX]
 +\gamma_{XX}\Dc[XX]
 +\gamma_{\mathrm{col}}\Dc[S_x],
 \qquad S_x=IX+XI
 }
 \label{eq:main_generator}
\end{equation}
である．各係数の物理的起源は
\begin{align}
 h_{XX}
 &:\quad \text{Fock依存$XX$回転角の熱平均},\nonumber\\
 \gamma_{XX}
 &:\quad \text{Fock依存$XX$回転角の熱分散},\nonumber\\
 \gamma_{\mathrm{col}}
 &:\quad \text{残留スピン振動エンタングルメント}
 \label{eq:physical_origins}
\end{align}
である．

\section{有限LD演算子とFock依存力}

矩形パルスに対する相互作用Hamiltonianを
\begin{equation}
 H(t)=A S_x
 \left(
 B e^{-\ii\delta t}+B^\dagger e^{\ii\delta t}
 \right)
 \label{eq:hamiltonian}
\end{equation}
とする．現在の数値実装で用いている運動演算子は
\begin{equation}
 B
 =a-\frac{\eta^2}{2}(N+1)a
 =a f(N),
 \qquad
 f(N)=1-\frac{\eta^2}{2}N
 \label{eq:B_operator}
\end{equation}
である．したがってFock状態に対して
\begin{equation}
 B\ket{n}
 =\sqrt{n}\,f_n\ket{n-1},
 \qquad
 f_n=1-\frac{\eta^2}{2}n
 \label{eq:B_on_fock}
\end{equation}
と作用する．通常のLD近似では$f_n=1$であり，すべてのFock状態が
同じスピン依存力を受ける．有限LD補正を含めると力の強さが$n$に
依存し，Fock状態ごとに異なる幾何学的位相と残留軌道を生じる．

\section{第二Magnus項とFock依存位相}

\subsection{交換子のFock表示}

$S_x$の固有値を$s\in\{2,0,0,-2\}$とする．各スピン枝での
Hamiltonianは
\begin{equation}
 H_s(t)
 =g_s\left(
 B e^{-\ii\delta t}+B^\dagger e^{\ii\delta t}
 \right),
 \qquad g_s=sA
 \label{eq:branch_hamiltonian}
\end{equation}
である．一周期$T=2\pi/\delta$では第一Magnus項は消え，第二項は
\begin{equation}
 \Omega_{2,s}
 =\ii\,2\pi
 \left(\frac{g_s}{\delta}\right)^2
 [B,B^\dagger]
 \label{eq:omega2}
\end{equation}
となる．交換子はFock基底で対角であり，
\begin{equation}
 [B,B^\dagger]\ket{n}
 =\Delta_n\ket{n}
\end{equation}
を満たす．式\eqref{eq:B_on_fock}から
\begin{align}
 \Delta_n
 &=(n+1)f_{n+1}^2-nf_n^2\nonumber\\
 &=1-\eta^2(2n+1)
 +\frac{\eta^4}{4}(3n^2+3n+1)
 \label{eq:delta_n}
\end{align}
を得る．

\subsection{$XX$回転角}

$s=2$のforced branchと$s=0$のnull branchとの相対位相を
$2\theta_n$と定義する．式\eqref{eq:omega2}より
\begin{equation}
 \boxed{
 \theta_n^{(2)}
 =\theta_{\mathrm{LD}}\Delta_n,
 \qquad
 \theta_{\mathrm{LD}}
 =4\pi\left(\frac{A}{\delta}\right)^2
 }
 \label{eq:theta_n}
\end{equation}
となる．現在のパラメータ
\begin{equation}
 A=0.125,\qquad
 \delta=0.5,\qquad
 \eta=0.1
\end{equation}
では$\theta_{\mathrm{LD}}=\pi/4$であり，
\begin{equation}
 \theta_n^{(2)}
 =0.777564
 -0.0156490\,n
 +5.89049\times10^{-5}n^2
 \label{eq:theta_n_numerical_analytic}
\end{equation}
を得る．一方，$n=0,\ldots,51$のfull numerical evolutionは
\begin{equation}
 \theta_n^{\mathrm{num}}
 \simeq
 0.771842
 -0.0154539\,n
 +5.75093\times10^{-5}n^2
 \label{eq:theta_n_fit}
\end{equation}
で近似でき，最大残差は$2.66\times10^{-5}\,\mathrm{rad}$である．

第二Magnus項は$n$依存性をよく説明するが，ほぼ$n$に依存しない
$-5.75\times10^{-3}\,\mathrm{rad}$程度のオフセットが残る．これは
強駆動および高次Magnus項による基準角のずれである．以下では$n=0$の
一点でこのオフセットを較正し，
\begin{equation}
 \theta_n^{\mathrm{ana}}
 =
 \theta_0^{\mathrm{num}}
 +\theta_{\mathrm{LD}}(\Delta_n-\Delta_0)
 \label{eq:one_point_calibration}
\end{equation}
とする．この較正後，温度依存部分には自由パラメータがない．

\section{熱平均チャネルと三成分への写像}

\subsection{スピン枝ごとのチャネル}

運動初期状態を平均フォノン数$\bar n=\mu$の熱状態
\begin{equation}
 \rho_{\mathrm{th}}
 =\sum_{n=0}^{\infty}p_n\ket{n}\bra{n},
 \qquad
 p_n=(1-q)q^n,
 \qquad
 q=\frac{\mu}{1+\mu}
 \label{eq:thermal_distribution}
\end{equation}
とする．$S_x$固有空間への射影を$\Pi_s$とすれば，一周期の全時間発展は
\begin{equation}
 U(T)=\sum_s\Pi_s\otimes U_s(T)
\end{equation}
と分解できる．運動をtrace outしたスピンチャネルは
\begin{equation}
 \mathcal{E}
 \left(\Pi_s\rho\Pi_{s'}\right)
 =
 G_{ss'}\Pi_s\rho\Pi_{s'},
 \qquad
 G_{ss'}
 =\tr\left[
 U_s\rho_{\mathrm{th}}U_{s'}^\dagger
 \right]
 \label{eq:branch_channel}
\end{equation}
である．したがって，各スピン枝間のcoherenceは運動波束のoverlap
$G_{ss'}$によって決まる．

forced--null overlapをFock状態ごとに
\begin{equation}
 c_n
 =\bra{n}U_{s=2}(T)\ket{n}
 =r_ne^{2\ii\theta_n}
 \label{eq:c_n}
\end{equation}
と定義する．位相のみの熱平均と，Fock leakageも含む完全なoverlapは
\begin{align}
 Z&=\sum_{n=0}^{\infty}p_ne^{2\ii\theta_n},
 \label{eq:Z_definition}\\
 C&=\sum_{n=0}^{\infty}p_nc_n
 =\sum_{n=0}^{\infty}p_nr_ne^{2\ii\theta_n}
 \label{eq:C_definition}
\end{align}
である．

\subsection{三係数の抽出}

目標ゲート角を$\theta_\star=\pi/4$とする．理想ゲートを除いた
forced--null coherenceに式\eqref{eq:main_generator}を作用させると，
\begin{equation}
 \exp(-2\ii h_{XX})
 \exp(-2\gamma_{XX})
 \exp(-2\gamma_{\mathrm{col}})
 \label{eq:generator_coherence}
\end{equation}
を得る．これを$e^{-2\ii\theta_\star}C$と比較すれば，
\begin{align}
 \boxed{
 h_{XX}
 =\theta_\star-\frac{1}{2}\arg C
 }
 \label{eq:hxx_exact}\\
 \boxed{
 \gamma_{XX}
 =-\frac{1}{2}\log|Z|
 }
 \label{eq:gammaxx_exact}\\
 \boxed{
 \gamma_{\mathrm{col}}
 =-\frac{1}{2}\log|C|-\gamma_{XX}
 }
 \label{eq:gammacol_exact}
\end{align}
を得る．これらはforced--null coherenceを厳密に再現する．
$r_n$と$\theta_n$の相関が弱い場合，式\eqref{eq:hxx_exact}の$C$を
$Z$で置き換えてもほぼ同じ結果となる．

\section{平均角ずれと位相分散の閉形式}

\subsection{Fock数に対する線形近似}

Fock依存角を
\begin{equation}
 \theta_n=\theta_0-\kappa n
\end{equation}
と近似すると，式\eqref{eq:Z_definition}は幾何級数となり，
\begin{equation}
 Z
 =e^{2\ii\theta_0}
 \frac{1-q}{1-qe^{-2\ii\kappa}}
 \label{eq:Z_linear}
\end{equation}
と閉じた形で書ける．したがって
\begin{align}
 \theta_{\mathrm{eff}}
 &\equiv\frac{1}{2}\arg Z\nonumber\\
 &=\theta_0-\frac{1}{2}
 \operatorname{atan2}
 \left(
 q\sin2\kappa,\,
 1-q\cos2\kappa
 \right),
 \label{eq:theta_eff_linear}\\
 \gamma_{XX}
 &=\frac{1}{4}\log
 \left[
 \frac{1-2q\cos2\kappa+q^2}{(1-q)^2}
 \right].
 \label{eq:gammaxx_linear}
\end{align}
特に$|\kappa|\ll1$では
\begin{equation}
 \boxed{
 h_{XX}
 \simeq
 (\theta_\star-\theta_0)+\kappa\mu,
 \qquad
 \gamma_{XX}
 \simeq
 \kappa^2\mu(1+\mu)
 }
 \label{eq:linear_small_kappa}
\end{equation}
となる．平均角ずれはほぼ$\bar n$に比例し，位相分散は熱分布の分散
$\bar n(1+\bar n)$に比例する．

\subsection{二次のFock依存性}

式\eqref{eq:delta_n}に対応して
\begin{equation}
 \theta_n=\theta_0+bn+cn^2
\end{equation}
まで残す．cumulant展開の最低次では
\begin{equation}
 \theta_{\mathrm{eff}}\simeq\E[\theta_n],
 \qquad
 \gamma_{XX}\simeq\Var(\theta_n)
 \label{eq:cumulant_leading}
\end{equation}
である．熱分布のモーメントから
\begin{equation}
 \E[\theta_n]
 =\theta_0+b\mu+c(2\mu^2+\mu)
 \label{eq:theta_mean_quadratic}
\end{equation}
および
\begin{align}
 \Var(\theta_n)
 ={}&
 b^2\mu(1+\mu)\nonumber\\
 &+2bc\,\mu(1+\mu)(4\mu+1)\nonumber\\
 &+c^2\mu(1+\mu)(20\mu^2+12\mu+1)
 \label{eq:theta_variance_quadratic}
\end{align}
を得る．より高精度には，式\eqref{eq:one_point_calibration}を
式\eqref{eq:Z_definition}の幾何重み付き級数に直接代入すればよい．

\section{第三Magnus項と残留運動リーク}

有限LD演算子では$[B,B^\dagger]$が定数でないため，高次交換子が
消えない．一周期の第三Magnus項は
\begin{equation}
 \Omega_{3,s}
 =\ii\,2\pi
 \left(\frac{g_s}{\delta}\right)^3Q
 \label{eq:omega3}
\end{equation}
であり，
\begin{equation}
 Q
 =[B,[B,B^\dagger]]
 +[B^\dagger,[B^\dagger,B]]
 \label{eq:Q_operator}
\end{equation}
と定義する．式\eqref{eq:delta_n}を用いると
\begin{equation}
 Q\ket{n}
 =q_n\ket{n-1}+q_{n+1}\ket{n+1},
\end{equation}
\begin{equation}
 \boxed{
 q_n
 =\sqrt{n}\,f_n(\Delta_n-\Delta_{n-1}),
 \qquad q_0=0
 }
 \label{eq:q_n}
\end{equation}
を得る．したがってforced branch
$(s=2,\ g_s=2A)$における弱リーク確率は
\begin{equation}
 \boxed{
 \ell_n^{(3)}
 \simeq
 \left[
 2\pi
 \left(\frac{2A}{\delta}\right)^3
 \right]^2
 \left(q_n^2+q_{n+1}^2\right)
 }
 \label{eq:leakage_n}
\end{equation}
となる．$r_n=\sqrt{1-\ell_n}\simeq1-\ell_n/2$を
式\eqref{eq:C_definition}に代入し，位相分散を分離すると
\begin{equation}
 \boxed{
 \gamma_{\mathrm{col}}
 \simeq
 \frac{1}{4}\sum_{n=0}^{\infty}p_n\ell_n^{(3)}
 }
 \label{eq:gammacol_leakage}
\end{equation}
を得る．

この減衰は$s$の異なるスピン枝が異なる残留運動状態を持つために生じる．
したがって結合演算子は$S_x=IX+XI$である．実際，
\begin{equation}
 \Dc[IX+XI]
 =\Dc[IX]+\Dc[XI]+\Cc_{IX,XI}
 \label{eq:collective_decomposition}
\end{equation}
であり，ここで
\begin{equation}
 \Cc_{P,Q}(\rho)
 =
 P\rho Q+Q\rho P
 -\frac{1}{2}\{PQ+QP,\rho\}.
\end{equation}
よって，数値的なgeneratorにおいて$\Dc[IX]$，$\Dc[XI]$，
$\Cc_{IX,XI}$の係数がほぼ等しいことも説明できる．

先頭の$O(\eta^2)$だけを残すと
\begin{equation}
 q_n\simeq-2\eta^2\sqrt{n}
\end{equation}
であるため，
\begin{equation}
 \gamma_{\mathrm{col}}
 \propto
 \eta^4
 \left(\frac{2A}{\delta}\right)^6
 (2\bar n+1)
 \label{eq:gammacol_scaling}
\end{equation}
という温度・パラメータ依存性を得る．定量比較には
式\eqref{eq:q_n}の$f_n,\Delta_n,q_n$を用いた熱和を使用する．

\section{数値結果との比較}

\subsection{平均角ずれと位相分散}

$n=0$の位相だけを較正し，第二Magnus項から得た$\Delta_n$を用いて
温度依存性を予測した結果を表\ref{tab:phase_comparison}に示す．

\begin{table}[htbp]
 \centering
 \small
 \caption{平均角ずれと位相分散の比較．}
 \label{tab:phase_comparison}
 \setlength{\tabcolsep}{5pt}
 \begin{tabular}{c@{\quad}rrrr}
  \toprule
  $\bar n$
  & $h_{XX}^{\mathrm{num}}$
  & $h_{XX}^{\mathrm{ana}}$
  & $\gamma_{XX}^{\mathrm{num}}$
  & $\gamma_{XX}^{\mathrm{ana}}$\\
  \midrule
  $0.01$ & $0.0137369$ & $0.0137390$ & $2.3913{\times}10^{-6}$ & $2.4541{\times}10^{-6}$\\
  $1$    & $0.0288390$ & $0.0290554$ & $4.5936{\times}10^{-4}$ & $4.7158{\times}10^{-4}$\\
  $2$    & $0.0438260$ & $0.0442921$ & $1.3352{\times}10^{-3}$ & $1.3720{\times}10^{-3}$\\
  $3$    & $0.0585297$ & $0.0592932$ & $2.5850{\times}10^{-3}$ & $2.6602{\times}10^{-3}$\\
  $4$    & $0.0729375$ & $0.0740587$ & $4.1656{\times}10^{-3}$ & $4.2966{\times}10^{-3}$\\
  \bottomrule
 \end{tabular}
\end{table}

$\bar n=4$においても，平均角ずれの相対差は約$1.5\%$，
位相分散の相対差は約$3.1\%$である．

\subsection{集団$X$成分}

第三Magnus項の$q_n$を熱平均した予測を
表\ref{tab:collective_comparison}に示す．

\begin{table}[htbp]
 \centering
 \small
 \caption{集団$X$成分の比較．}
 \label{tab:collective_comparison}
 \setlength{\tabcolsep}{8pt}
 \begin{tabular}{crrr}
  \toprule
  $\bar n$
  & $\gamma_{\mathrm{col}}^{\mathrm{num}}$
  & $\gamma_{\mathrm{col}}^{\Omega_3}$
  & 相対差\\
  \midrule
  $0.01$ & $5.9979{\times}10^{-5}$ & $6.1330{\times}10^{-5}$ & $+2.25\%$\\
  $1$    & $1.6774{\times}10^{-4}$ & $1.7172{\times}10^{-4}$ & $+2.37\%$\\
  $2$    & $2.6556{\times}10^{-4}$ & $2.7241{\times}10^{-4}$ & $+2.58\%$\\
  $3$    & $3.5315{\times}10^{-4}$ & $3.6312{\times}10^{-4}$ & $+2.82\%$\\
  $4$    & $4.3146{\times}10^{-4}$ & $4.4470{\times}10^{-4}$ & $+3.07\%$\\
  \bottomrule
 \end{tabular}
\end{table}

さらに，full numerical evolutionから直接得たFock leakage
\begin{equation}
 \ell_n=1-|c_n|^2
\end{equation}
を式\eqref{eq:gammacol_leakage}に代入すると，$\bar n\leq4$の
全範囲で$\gamma_{\mathrm{col}}$と$0.84\%$以内で一致する．したがって
約$3\%$の残差は主として第三Magnus近似の打ち切りに由来し，
$\gamma_{\mathrm{col}}$をFock leakageの熱平均と解釈すること自体は
高精度に成立する．

\section{物理的帰結}

以上から，数値的に得られた三成分は独立なfit parameterではなく，
同じ有限LD Hamiltonianが作るFock分解の異なる統計量として
\begin{equation}
 \boxed{
 \begin{aligned}
  h_{XX}
  &\longleftarrow
  \text{Fock依存位相の平均},\\
  \gamma_{XX}
  &\longleftarrow
  \text{Fock依存位相の分散},\\
  \gamma_{\mathrm{col}}
  &\longleftarrow
  \text{Fock依存の残留軌道リーク}
 \end{aligned}
 }
 \label{eq:unified_picture}
\end{equation}
と統一的に理解できる．

単一のパルス面積またはdetuning較正で除去できるのは平均角ずれ
$h_{XX}$だけである．$\gamma_{XX}$と$\gamma_{\mathrm{col}}$は
熱分布による混合および残留スピン運動絡み合いに由来するため，
単一のユニタリ較正後にも誤差床として残る．したがって，
「平均MS回転角を較正しても有限温度誤差は消えず，確率的$XX$成分と
集団$X$成分へ分離される」という結論が得られる．

\section{適用範囲と残された課題}

本解析で使用した$B$は，第一側帯波演算子を$O(\eta^2)$まで保持した
Hamiltonianである．そのHamiltonian自体は数値的に厳密時間発展
させているが，光学相互作用の全次数表式ではない．また，現時点では
次の二点を解析的に完全には取り込んでいない．
\begin{enumerate}
 \item $n$にほぼ依存しない強駆動・高次Magnus由来の回転角オフセット．
 \item 第三Magnus項より高次の小さなFock leakage補正．
\end{enumerate}
したがって本結果は「全次数の厳密解」ではなく，
\begin{quote}
 $n=0$の一点較正を施した第二Magnus予測と第三Magnus予測が，
 full numerical evolutionおよびQPTから得た三成分を
 $\bar n\leq4$で約$3.1\%$以内に説明する
\end{quote}
と位置付けるのが適切である．

\section*{対応する実装と数値データ}

\begin{itemize}
 \item \texttt{ms\_gate\_functions.py}: 有限LD Hamiltonian．
 \item \texttt{phonon\_xx\_angle\_analysis.py}:
       $c_n$, $Z$, $C$と三係数の抽出．
 \item \texttt{noise\_error\_structure.py}:
       QPT generatorの基底規約．
 \item \texttt{results/noise\_rate\_nbar\_sweep/noise\_free\_diagnostics/}\\
       \texttt{fock\_resolved/fock\_resolved\_curve.csv}:
       Fock分解データ．
 \item \texttt{results/noise\_rate\_nbar\_sweep/noise\_free\_diagnostics/}\\
       \texttt{fock\_resolved/fock\_thermal\_summary.csv}:
       熱平均結果．
\end{itemize}

\end{document}
